\documentclass[11pt]{article}
\usepackage{mathblatt}
% gesetzt von werkzeuge/setzer.py; Übersicht nach bau/bauregeln.md „Übersicht und Serie“.
\newcommand{\kennung}{T74}
\newcommand{\bereich}[1]{\par\addvspace{14pt}\noindent{\large\bfseries #1}\par\nobreak\vspace{4pt}}
\newcommand{\bl}[1]{\par\noindent #1\par\vspace{3pt}}
\newcommand{\blrand}[1]{\par\noindent{\color{mbgrau}#1}\par\vspace{3pt}}
\newcommand{\blhier}[1]{\par\noindent\llap{$\blacktriangleright$\hspace{0.5em}}\textbf{#1}\par\vspace{3pt}}
\begin{document}
\pfheftstil
\fancyfoot[C]{{\footnotesize\color{mbgrau}\kennung\hfill\thepage}}
\pfheftkopf{Daten}{Klasse 6}

\bereich{Vorher}
\bl{Zahlen ordnen und Skalen ablesen}

\bereich{Daten darstellen}
\bl{Häufigkeiten}
\blhier{Säulen-, Balken- und Liniendiagramme}
\bl{Streifen- und Kreisdiagramm}

\bereich{Daten auswerten}
\bl{Kenngrößen}
\bl{Diagramme beurteilen und Boxplot}

\bereich{Weiter}
\bl{Vierfeldertafel}
\end{document}
